\subsection{Stanley-Reisner rings}

Let $\Delta$ be a simplicial complex of dimension $n-1$  with vertices $v_1, \ldots, v_N$. The Stanley-Reisner ring of  $\Delta$ over a field $K$ is
\[ \cA(\Delta) = K[x_1,\ldots, x_N]/I_\Delta,\]
where $I_\Delta$ is the ideal generated by all square-free monomials $\prod_{i\in S} x_i$ such that the set $\{ v_i \}_{i\in S}$ is not a simplex in $\Delta$. The ring $\cA(\Delta)$ is a graded $K$-algebra. Given homogeneous degree $1$ elements $\theta_1,\ldots, \theta_n \in \cA^1(\Delta)$, we define the cohomology ring
\[ H(\Delta) = \cA(\Delta)/(\theta_1,\ldots,\theta_n).\]
To remove dependence on the choice of $\theta_i$, we work with generic parameters
\[ \theta_i = \a{i, 1} x_1 + \a{i, 2} x_2 + \cdots + \a{i, N} x_N, \quad i=1,\ldots,n,\]
where $\a{i, j}$ are indeterminates and the field $K$ is the field of rational functions $K=k(\a{i, j})$. We will only consider this generic case. 

If $\Delta$ is a homology sphere over $k$ and the parameters are generic as above, then the ring $H(\Delta)$ is a standard graded, Artinian, Gorenstein $K$-algebra of socle degree $n$. The \Po pairing defined by multiplication 
\[ H^m(\Delta) \times H^{n-m}(\Delta) \longrightarrow H^n(\Delta) \isom K\]
is a nondegenerate  bilinear pairing. If $\Delta$ is only an oriented pseudo-manifold over $k$, then we still have $H^n(\Delta) \isom K$, but the pairing may be degenerate.
